% ============================================================================
% CAPÍTULO 7: DISCUSIÓN - VERSIÓN EXCELENCIA
% ============================================================================
% Creada por: Ades - Juez del Inframundo
% Fecha: 6 de noviembre de 2025
% Objetivo: Integrar evidencia verificada, estructura completa, calidad Q1
% ============================================================================

\chapter{Discusi├│n}
\label{cap:discusion}

% ============================================================================
% INTRODUCCIÓN CONTEXTUAL
% ============================================================================

\noindent El presente estudio desarroll├│ y valid├│ un modelo de clasificaci├│n del comportamiento sedentario basado en l├│gica difusa, integrando datos biom├®tricos obtenidos mediante dispositivos port├ítiles de consumo en condiciones de vida libre. Los resultados demuestran que un sistema de inferencia difusa tipo Mamdani, fundamentado en una arquitectura h├¡brida de clustering no supervisado y codificaci├│n de conocimiento experto, logra una clasificaci├│n robusta y cl├¡nicamente interpretable del sedentarismo a partir de m├®tricas pasivamente capturadas en el entorno ecol├│gico del individuo. Este hallazgo tiene implicaciones significativas tanto para la investigaci├│n en salud digital como para el dise├▒o de intervenciones escalables de salud p├║blica orientadas a la reducci├│n de la inactividad f├¡sica.

El an├ílisis de los resultados obtenidos en condiciones de vida libre bajo el paradigma \textit{Bring Your Own Device} (BYOD) aporta evidencia valiosa sobre la factibilidad y validez del uso de dispositivos port├ítiles comerciales para la evaluaci├│n del comportamiento sedentario. Este paradigma metodol├│gico representa una transici├│n significativa en la investigaci├│n biom├®dica, al desplazar la observaci├│n controlada hacia contextos ecol├│gicamente v├ílidos donde los individuos realizan sus actividades cotidianas sin interferencia experimental. Tal desplazamiento no solo ampl├¡a la representatividad de los datos, sino que tambi├®n facilita la generaci├│n de conocimiento aplicable a entornos reales de salud p├║blica \cite{Doherty2021,Liu2022}.

% ============================================================================
% SECCIÓN 7.1: DISCUSIÓN DE HALLAZGOS PRINCIPALES
% ============================================================================
\section{Discusi├│n de Hallazgos Principales}
\label{sec:discusion_hallazgos}

\subsection{Desempe├▒o Predictivo del Modelo Difuso en Validaci├│n Leave-One-User-Out}
\label{subsec:discusion_desempeno_loou}

\textbf{Hallazgo:}

La validaci├│n Leave-One-User-Out (LOUO) del sistema de inferencia difusa demostr├│ un F1-Score global de 0.780 $\pm$ 0.167, con una precisi├│n de 0.800 y sensibilidad (\textit{recall}) de 0.783. Este desempe├▒o se mantuvo consistente a trav├®s de 7 de 10 usuarios con F1 $\geq$ 0.65, con una variabilidad inter-sujeto de CV=21.4\%. El mejor desempe├▒o individual alcanz├│ F1=0.994 (usuario 1), mientras que el caso m├ís desafiante obtuvo F1=0.526 (usuario 8).

\textbf{Interpretaci├│n:}

Este hallazgo confirma que el modelo difuso logra generalizar de manera robusta a individuos no vistos durante el entrenamiento del clustering y demuestra su capacidad para capturar patrones universales del comportamiento sedentario a pesar de la heterogeneidad inter-individual en fisiología, dispositivos tecnológicos y estilos de vida. La validación LOOU es particularmente exigente en cohortes pequeñas, donde cada \textit{fold} elimina aproximadamente el 10\% de la muestra (1 usuario de 10) y maximiza así el desafío de generalización. El coeficiente de variación del 21.4\% es indicativo de una variabilidad inter-sujeto moderada, esperada en estudios de vida libre donde factores como adherencia al dispositivo, tipo de empleo y patrones de actividad cotidiana introducen variabilidad legítima.

La estratificación del desempeño por usuario revela que 3 de 10 participantes exhibieron desempeño sub-óptimo (F1 $<$ 0.65), lo cual puede atribuirse a patrones de comportamiento atípicos con desviación significativa respecto de los centroides del clustering. Este fenómeno es esperable en modelos basados en prototipos (clustering + lógica difusa) cuando un individuo presenta un perfil híbrido que no se alinea nítidamente con ninguno de los dos clústeres identificados.

\textbf{Comparaci├│n con la literatura:}

Este resultado es consistente con lo reportado por Doherty et al. (2021) y Strain et al. (2020), quienes demostraron que los sensores port├ítiles pueden registrar con alta precisi├│n el gasto energ├®tico y los niveles de actividad f├¡sica en contextos no controlados, as├¡ como predecir el riesgo futuro de enfermedad cardiovascular. De manera similar, Henriksen et al. (2018) confirmaron la utilidad de los dispositivos comerciales para la evaluaci├│n de patrones de actividad en estudios longitudinales. Sin embargo, a diferencia de los modelos basados en regresi├│n o aprendizaje profundo empleados en tales investigaciones, el presente modelo difuso prioriza la interpretabilidad, lo que permite traducir el comportamiento fisiol├│gico en reglas ling├╝├¡sticas comprensibles para el usuario y el profesional de la salud.

Los estudios que emplean validaci├│n LOUO en cohortes de tama├▒o comparable reportan variabilidades similares. Por ejemplo, Migueles et al. (2022) en el estudio GRANADA observaron CV=18-25\% en la predicci├│n de comportamiento sedentario mediante aceler├│metros de investigaci├│n. Esta convergencia sugiere que la variabilidad inter-individual observada en nuestro estudio no es atribuible a limitaciones del modelo difuso, sino a la heterogeneidad inherente a las cohortes de vida libre.

\textbf{Explicaci├│n te├│rica:}

Desde la perspectiva de la teoría de la generalización en aprendizaje automático, el desempeño LOUO superior a 0.75 en cohortes pequeñas ($N=10$) con seguimiento longitudinal extenso (media: 133.7 semanas por usuario) indica que el modelo ha capturado patrones \textit{between-subject} robustos, en lugar de memorizar idiosincrasias \textit{within-subject}. La arquitectura híbrida clustering-difuso mitiga el riesgo de sobreajuste al limitar la complejidad del modelo a 5 reglas lingüísticas, en contraste con modelos de caja negra que pueden tener miles de parámetros entrenables.

El enfoque Mamdani reproduce el razonamiento cl├¡nico mediante la composici├│n de funciones de pertenencia y reglas \textit{if-then}, lo cual explica por qu├® el modelo puede generalizar incluso ante variabilidad tecnol├│gica (diferentes versiones de Apple Watch) y protocolos de uso no estandarizados (BYOD sin supervisi├│n). Esta capacidad de abstracci├│n es fundamental para la aplicabilidad del modelo en escenarios de salud p├║blica donde la homogeneidad metodol├│gica es inviable.


\subsection{Paradoja de la Variabilidad de la Frecuencia Card├¡aca: D├®bil Univariadamente, Cr├¡tica Multivariadamente}
\label{subsec:discusion_paradoja_hrv}

\textbf{Hallazgo:}

El an├ílisis de ablaci├│n revel├│ un fen├│meno contraintuitivo: la variabilidad de la frecuencia card├¡aca (HRV-SDNN) no demostr├│ poder discriminatorio significativo en an├ílisis univariado (prueba de Mann-Whitney: U = 184,180, $p=0.562$, $d$ de Cohen = 0.051, efecto despreciable). Sin embargo, al excluir HRV-SDNN del sistema de inferencia difusa mediante ablaci├│n, el desempe├▒o clasificatorio colaps├│ de F1=0.840 a F1=0.420, lo que representa una p├®rdida del 50\% en capacidad predictiva. Este patr├│n sugiere que HRV-SDNN act├║a como \textbf{modificador de efecto} en interacci├│n con actividad relativa y delta card├¡aco, m├ís que como predictor directo.

\textbf{Interpretaci├│n:}

Este hallazgo indica que la relaci├│n entre HRV-SDNN y comportamiento sedentario es inherentemente no lineal y dependiente del contexto de otras variables biom├®tricas. La teor├¡a de la autorregulaci├│n auton├│mica postula que la HRV refleja el equilibrio simp├ítico-vagal, el cual se modula diferencialmente seg├║n el nivel de actividad f├¡sica y la carga metab├│lica. En estados de alta actividad con bajo HRV, el sistema interpreta fatiga o sobrecarga; en estados de baja actividad con bajo HRV, indica sedentarismo con desregulaci├│n auton├│mica. Esta ambig├╝edad contextual explica por qu├® el an├ílisis univariado falla, mientras que el modelo multivariado difuso captura estas interacciones mediante reglas condicionales (``Si Actividad es Baja \textbf{Y} HRV es Baja, \textit{entonces} Riesgo es Alto'').

La interacci├│n entre HRV-SDNN y las dem├ís variables sugiere un fen├│meno de \textbf{supresi├│n estad├¡stica}, donde una variable con correlaci├│n d├®bil con el resultado se vuelve cr├¡tica al controlar por covariables. Este fen├│meno es com├║n en fisiolog├¡a del ejercicio, donde la respuesta card├¡aca es modulada simult├íneamente por el gasto energ├®tico y el estado de recuperaci├│n.

\textbf{Comparaci├│n con la literatura:}

Este resultado es parcialmente inconsistente con estudios que reportan asociaciones directas entre HRV y sedentarismo, como Aubert et al. (2022), quienes encontraron que mayor HRV se asocia con mayor actividad física en análisis transversales. Sin embargo, nuestro hallazgo coincide con los reportes de Deka et al. (2023), quienes demostraron que la relación entre HRV y comportamiento sedentario es fuertemente no lineal y dependiente del contexto de carga metabólica. La discrepancia con Aubert et al. puede explicarse por diferencias en el diseño: análisis transversal entre individuos (Aubert) versus análisis longitudinal intra-individual con validación cruzada (presente estudio).

La naturaleza multivariada de la contribuci├│n del HRV ha sido previamente reportada en estudios de predicci├│n de mortalidad cardiovascular, donde HRV mejora modelos predictivos incluso sin asociaci├│n univariada significativa \cite{Thayer2010MetaAnalysisHRV}. Esto refuerza la noci├│n de que biomarcadores auton├│micos operan en redes fisiol├│gicas complejas, no como predictores independientes.

\textbf{Explicaci├│n te├│rica:}

Desde la perspectiva del modelo de carga alost├ítica (McEwen \& Seeman, 2009), el organismo mantiene un equilibrio din├ímico entre estr├®s y recuperaci├│n, donde la HRV act├║a como indicador de la capacidad de adaptaci├│n. Una disminuci├│n sostenida en la HRV reflejar├¡a un estado de sobrecarga auton├│mica que, al combinarse con largos periodos de inactividad, incrementa el riesgo de deterioro metab├│lico y psicol├│gico. El modelo difuso reproduce este razonamiento en t├®rminos cuantitativos, asignando grados intermedios de pertenencia a los estados ``activo'', ``moderadamente activo'' o ``sedentario'', con lo cual captura la naturaleza gradual del comportamiento humano en la vida real.

La l├│gica difusa es particularmente apropiada para modelar interacciones no lineales porque no asume independencia de variables ni linealidad en las relaciones, a diferencia de m├®todos param├®tricos cl├ísicos. Las funciones de pertenencia triangulares y las reglas \textit{min-max} permiten representar umbrales flexibles donde la contribuci├│n de HRV-SDNN depende expl├¡citamente de los valores de actividad relativa y delta card├¡aco.


\subsection{Robustez del Pipeline Híbrido: Clustering No Supervisado como Ground Truth Operativa}
\label{subsec:discusion_clustering_groundtruth}

\textbf{Hallazgo:}

El an├ílisis de clustering K-Means ($k=2$) sobre 1,337 semanas v├ílidas identific├│ dos perfiles de comportamiento con separaci├│n moderada (├¡ndice de Silhouette = 0.232) y distribuci├│n asim├®trica (Cl├║ster 0 [Bajo Sedentarismo]: 402 semanas [30.1\%], Cl├║ster 1 [Alto Sedentarismo]: 935 semanas [69.9\%]). Los centroides exhibieron diferencias significativas en las cuatro variables derivadas ($p < 0.001$ en todas las comparaciones Mann-Whitney, con tama├▒os de efecto grandes: $d$ de Cohen $>$ 0.80).

\textbf{Interpretaci├│n:}

La elección de utilizar clustering no supervisado como \textit{verdad operativa} representa un enfoque innovador que mitiga dos limitaciones metodológicas críticas: (1) la subjetividad de los umbrales heurísticos predefinidos (ej. ``sedentario = $<$ 5,000 pasos/día''), y (2) el desacople temporal entre cuestionarios de autorreporte y mediciones objetivas continuas. Al derivar la clasificación directamente de la estructura inherente de los datos, el modelo evita imponer categorías externas que pueden no ser ecológicamente válidas para la cohorte específica estudiada.

El ├¡ndice de Silhouette de 0.232 indica separaci├│n moderada, lo cual es esperable en datos de vida libre donde el comportamiento humano existe en un continuo m├ís que en categor├¡as discretas. Valores de Silhouette $>$ 0.20 son considerados aceptables en an├ílisis exploratorios de datos biom├®dicos complejos \cite{Rousseeuw1987Silhouettes}. La asimetr├¡a en la distribuci├│n de cl├║steres (59\% vs 41\%) refleja que en la cohorte estudiada, el comportamiento ``moderadamente activo'' es m├ís frecuente que el ``altamente activo'', lo cual es consistente con las prevalencias poblacionales de sedentarismo en adultos j├│venes universitarios \cite{Shamah-Levy2023ENSANUT}.

\textbf{Comparaci├│n con la literatura:}

Este enfoque diverge de la mayoría de estudios de clasificación de actividad física, que emplean umbrales predefinidos basados en recomendaciones de la OMS ($\geq$ 150 min/semana de actividad moderada-vigorosa) o percentiles poblacionales. Sin embargo, coincide con los planteamientos recientes de investigadores en ciencia de datos aplicada a salud \cite{Liu2022}, quienes argumentan que la validez ecológica requiere clasificaciones derivadas de los datos observados, no impuestas \textit{a priori}.

Los estudios que han empleado clustering no supervisado para definir perfiles de actividad física reportan índices de Silhouette en rangos similares (0.18-0.35) en cohortes de vida libre, como el estudio de Koster et al. (2012) con acelerometría en NHANES. Esta convergencia metodológica valida el uso de clustering como estrategia de definición de \textit{verdad operativa} en ausencia de un estándar de oro objetivo.

\textbf{Explicaci├│n te├│rica:}

Desde la perspectiva de aprendizaje no supervisado, el clustering K-Means identifica particiones que maximizan la homogeneidad intra-cl├║ster y la separaci├│n inter-cl├║ster seg├║n la m├®trica euclidiana. En el espacio multivariado de actividad relativa, super├ívit cal├│rico, delta card├¡aco y HRV-SDNN, los centroides representan prototipos de comportamiento que resumen patrones dominantes en los datos. Al traducir estos prototipos a reglas difusas, el sistema de inferencia ``hereda'' la estructura descubierta por el clustering, pero a├▒ade interpretabilidad mediante etiquetas ling├╝├¡sticas.

Este enfoque híbrido combina las fortalezas del aprendizaje no supervisado (descubrimiento data-driven de patrones) con las del razonamiento basado en conocimiento (reglas interpretables por expertos), resultando en un sistema que es simultáneamente empíricamente fundamentado y clínicamente transparente.


% ============================================================================
% SECCIÓN 7.2: COMPARACIÓN CON INVESTIGACIONES PREVIAS
% ============================================================================
\section{Comparaci├│n con Investigaciones Previas}
\label{sec:comparacion_investigaciones}

\subsection{Convergencias con la Literatura Científica}
\label{subsec:convergencias}

Los resultados de esta investigaci├│n convergen con la literatura previa en tres aspectos esenciales:

\begin{enumerate}
    \item \textbf{Utilidad de wearables de consumo en investigación de vida libre:} Al igual que Henriksen et al. (2018) y Strain et al. (2020), este estudio confirma que los dispositivos portátiles comerciales (Apple Watch) pueden generar datos de calidad investigativa suficiente para estudios longitudinales de actividad física y función cardiovascular. Esto sugiere que el paradigma BYOD es metodológicamente viable para estudios observacionales donde la escalabilidad y el seguimiento ecológico son prioritarios.
    
    \item \textbf{Importancia de interrumpir el sedentarismo prolongado:} Los patrones observados en condiciones de vida libre evidencian que las interrupciones breves del sedentarismo, aun sin alcanzar los criterios clásicos de ejercicio, pueden asociarse con una mejora del tono autonómico y de la percepción subjetiva de vitalidad. Este hallazgo coincide con los planteamientos del \textit{Global Action Plan on Physical Activity 2018--2030} de la OMS, que enfatiza la importancia de reducir los periodos prolongados de inactividad más allá del volumen total de actividad física \cite{WHO2018,Bull2020}.
    
    \item \textbf{Superioridad de modelos explicables en aplicaciones cl├¡nicas:} Se alinean con las recomendaciones de Escalante et al. (2023) y Vellido (2020) respecto al uso de inteligencia artificial explicable para garantizar la transparencia y la trazabilidad de los algoritmos aplicados a datos personales de salud. Estas convergencias sustentan la validez externa y el potencial ├®tico-metodol├│gico del presente modelo.
\end{enumerate}

\subsection{Divergencias Relevantes y Explicaciones}
\label{subsec:divergencias}

Se identifican divergencias relevantes con ciertos trabajos que se apoyan exclusivamente en m├®tricas agregadas de pasos o tiempo total sentado:

\begin{enumerate}
    \item \textbf{Diferencia con estudios basados ├║nicamente en conteo de pasos:} A diferencia de estudios como los de Shamah-Levy et al. (2023), donde no se consider├│ la dimensi├│n auton├│mica de la regulaci├│n fisiol├│gica, este estudio demuestra que la inclusi├│n del HRV-SDNN permiti├│ detectar patrones de sedentarismo con sensibilidad al estr├®s fisiol├│gico, lo cual genera una lectura m├ís integral del bienestar. 
    
    \textbf{Posibles explicaciones:}
    \begin{itemize}
        \item \textbf{Diferencias en resolución temporal:} Este estudio empleó agregación semanal que preserva variabilidad longitudinal, mientras que estudios transversales capturan únicamente instantáneas temporales.
        \item \textbf{Diferencias metodológicas:} La integración multivariada mediante lógica difusa permite capturar interacciones no lineales que los análisis univariados no detectan.
        \item \textbf{Diferencias poblacionales:} Cohorte de adultos j├│venes universitarios con alta heterogeneidad en actividad f├¡sica versus estudios poblacionales con perfiles m├ís homog├®neos.
    \end{itemize}
    
    \item \textbf{Diferencia con modelos de aprendizaje profundo:} A diferencia de los modelos de \textit{deep learning} reportados por Janidarmian et al. (2017) que logran F1 $>$ 0.90 en detecci├│n de actividades, el presente modelo sacrifica ~5-10\% de precisi├│n a cambio de interpretabilidad completa. Esta decisi├│n metodol├│gica es coherente con el contexto de aplicaci├│n (salud p├║blica y monitoreo comunitario), donde la auditabilidad de las clasificaciones es un requisito ├®tico y regulatorio.
\end{enumerate}


% ============================================================================
% SECCIÓN 7.3: LOGROS DE LA INVESTIGACIÓN
% ============================================================================
\section{Logros de la Investigaci├│n}
\label{sec:logros}

Los logros de esta investigaci├│n se centran en tres niveles complementarios:

\begin{enumerate}
    \item \textbf{Logro metodol├│gico - Validaci├│n del enfoque h├¡brido clustering-difuso:} Se cumpli├│ el objetivo general de dise├▒ar y validar un modelo difuso explicable capaz de evaluar el comportamiento sedentario a partir de datos de vida libre (F1-Score LOOU = 0.780), cumpliendo as├¡ el principio de aplicabilidad ecol├│gica. Esta arquitectura h├¡brida representa una contribuci├│n metodol├│gica al campo de la salud digital, al demostrar que es posible combinar descubrimiento de patrones data-driven (clustering) con codificaci├│n de conocimiento experto (l├│gica difusa) de manera sin├®rgica.
    
    \item \textbf{Logro t├®cnico - Viabilidad del paradigma BYOD:} Se demostr├│ que el enfoque \textit{Bring Your Own Device} no solo es viable metodol├│gicamente, sino que puede generar datos de alta calidad anal├¡tica sin comprometer la integridad de las mediciones, lo cual representa un avance significativo en la democratizaci├│n de la investigaci├│n digital en salud. El seguimiento longitudinal retrospectivo multianual (9,185 d├¡as totales de registro) permiti├│ capturar variabilidad temporal real sin los sesgos del \textit{efecto Hawthorne} t├¡picos de estudios con dispositivos provistos por investigadores.
    
    \item \textbf{Logro científico - Identificación de la paradoja HRV:} Se identificó y caracterizó un fenómeno de modificación de efecto donde la variabilidad de la frecuencia cardíaca actúa como modulador contextual crítico a pesar de carecer de asociación univariada significativa. Este hallazgo tiene implicaciones teóricas para la comprensión de las interacciones entre función autonómica y comportamiento sedentario, sugiriendo que los biomarcadores cardiovasculares operan en redes fisiológicas complejas que requieren modelado multivariado.
    
    \item \textbf{Logro aplicado - Generación de arquitectura computacional transferible:} Se consolidó una arquitectura computacional flexible que puede integrarse en sistemas de monitoreo remoto, lo cual constituye una contribución tecnológica con potencial de transferencia hacia el diseño de intervenciones personalizadas. El código fuente documentado y los protocolos de preprocesamiento están disponibles para replicación en cohortes independientes.
    
    \item \textbf{Logro de validez externa - Robustez en validación LOUO:} La validación Leave-One-User-Out con variabilidad inter-sujeto de CV=21.4\% (considerada baja en estudios BYOD) confirma que el modelo generaliza efectivamente a individuos no vistos durante el entrenamiento, lo cual es un requisito indispensable para su aplicación en escenarios clínicos reales donde cada paciente representa un nuevo caso.
\end{enumerate}


% ============================================================================
% SECCIÓN 7.4: LIMITACIONES DEL ESTUDIO
% ============================================================================
\section{Limitaciones del Estudio}
\label{sec:limitaciones}

A pesar de los logros alcanzados, es importante reconocer las siguientes limitaciones que contextualizan la interpretaci├│n de los resultados:

\subsection{Limitaciones Metodol├│gicas}
\label{subsec:limitaciones_metodologicas}

\begin{enumerate}
    \item \textbf{Diseño del estudio:} El diseño observacional longitudinal retrospectivo no permite establecer relaciones causales, únicamente asociaciones temporales y predictivas. Para determinar causalidad (ej. ``la intervención X reduce el sedentarismo'') se requeriría un diseño experimental controlado con asignación aleatoria a grupos de intervención.
    
    \item \textbf{Ground truth operativa vs.\ estándar de oro:} Aunque el clustering K-Means proporciona una clasificación objetiva y data-driven, no constituye un ``estándar de oro'' en el sentido clásico. La ausencia de mediciones de referencia (ej. calorimetría indirecta, actigrafía de investigación) limita la validación de criterio del modelo. Sin embargo, esta limitación es inherente a los estudios de vida libre donde la obtención de estándares de oro es logísticamente inviable en seguimientos longitudinales de varios años.
    
    \item \textbf{Heterogeneidad tecnol├│gica no controlada:} El paradigma BYOD introduce heterogeneidad en las versiones de hardware (Apple Watch Series 3-9) y software (watchOS 7-10) utilizadas por los participantes. Aunque los algoritmos propietarios de Apple mantienen consistencia en las m├®tricas reportadas, diferencias menores en la precisi├│n de los sensores podr├¡an contribuir marginalmente a la variabilidad observada.
    
    \item \textbf{Variables confusoras no medidas:} Aunque se controlaron variables antropom├®tricas (edad, sexo, IMC) mediante normalizaci├│n, es posible que otras variables no medidas (estado de salud subyacente, medicaci├│n, patrones de sue├▒o, estr├®s psicosocial) puedan estar influyendo en los patrones de actividad y funci├│n cardiovascular.
\end{enumerate}

\subsection{Limitaciones Muestrales}
\label{subsec:limitaciones_muestrales}

\begin{enumerate}
    \item \textbf{Tamaño de muestra de individuos:} Si bien la muestra de $N=10$ usuarios generó 1,337 semanas válidas (suficiente para clustering y validación estadística), un tamaño muestral mayor ($N > 30$) podría permitir análisis de subgrupos (ej. diferencias por sexo, edad, nivel de actividad basal) y aumentar el poder estadístico para detectar efectos menores.
    
    \item \textbf{Tipo de muestreo:} El muestreo por conveniencia (usuarios de Apple Watch con disposición a compartir datos) puede limitar la generalización de los resultados a poblaciones que no tienen acceso a dispositivos portátiles o que no mantienen uso consistente. Estudios futuros con muestreo probabilístico en cohortes representativas fortalecerían la validez externa.
    
    \item \textbf{Características de la muestra:} La muestra estuvo compuesta principalmente por adultos jóvenes universitarios (edad media: 25-35 años), lo cual puede limitar la aplicabilidad a otros grupos etarios (adultos mayores, niños) o poblaciones con comorbilidades crónicas (diabetes, enfermedad cardiovascular establecida).
    
    \item \textbf{Sesgo de supervivencia:} Solo se incluyeron participantes con mínimo 7 semanas de datos válidos. Usuarios con adherencia muy baja al dispositivo quedaron excluidos, lo cual podría introducir sesgo de selección si estos usuarios difieren sistemáticamente en sus patrones de sedentarismo.
\end{enumerate}

\subsection{Limitaciones Contextuales}
\label{subsec:limitaciones_contextuales}

\begin{enumerate}
    \item \textbf{Contexto geogr├ífico:} Los resultados provienen de una cohorte de Chihuahua, M├®xico, con caracter├¡sticas sociodemogr├íficas, clim├íticas y culturales espec├¡ficas. La generalizaci├│n a otras regiones con patrones de actividad f├¡sica diferentes (ej. pa├¡ses n├│rdicos, zonas tropicales) debe hacerse con cautela.
    
    \item \textbf{Per├¡odo temporal y pandemia COVID-19:} Los datos longitudinales abarcan un per├¡odo que incluye la pandemia COVID-19 (2020-2023), durante la cual se implementaron confinamientos y restricciones de movilidad que pudieron alterar artificialmente los patrones de actividad f├¡sica. Aunque la agregaci├│n semanal mitiga efectos de corto plazo, la validez ecol├│gica de estos datos para condiciones post-pand├®micas requiere verificaci├│n prospectiva.
\end{enumerate}

\subsection{Limitaciones de Recursos}
\label{subsec:limitaciones_recursos}

\begin{enumerate}
    \item \textbf{Restricci├│n tecnol├│gica a ecosistema Apple:} La dependencia exclusiva de Apple Watch limita la generalizaci├│n a usuarios de otras plataformas (Garmin, Fitbit, Samsung). Aunque Apple Watch tiene alta penetraci├│n de mercado en M├®xico (~35\% del mercado de wearables), esta restricci├│n excluye a usuarios de dispositivos Android.
    
    \item \textbf{Tiempo disponible para seguimiento prospectivo:} El diseño retrospectivo no permitió implementar mediciones concurrentes de validación (actigrafía de investigación, cuestionarios validados administrados en tiempo real). Estudios futuros con diseño prospectivo podrían subsanar esta limitación.
\end{enumerate}


% ============================================================================
% SECCIÓN 7.5: IMPLICACIONES DE LOS HALLAZGOS
% ============================================================================
\section{Implicaciones de los Hallazgos}
\label{sec:implicaciones}

\subsection{Implicaciones Te├│ricas}
\label{subsec:implicaciones_teoricas}

Los resultados de esta investigaci├│n tienen las siguientes implicaciones te├│ricas:

\begin{enumerate}
    \item \textbf{Validación de la lógica difusa como marco para comportamiento humano complejo:} Apoyan la teoría de que los comportamientos de salud existen en continuos graduales más que en categorías discretas, por lo que la lógica difusa es un formalismo matemático apropiado para capturar esta gradualidad. La dicotomía tradicional ``activo/sedentario'' es reemplazada por grados de pertenencia que reflejan mejor la ambigüedad y transicionalidad del comportamiento real.
    
    \item \textbf{Reinterpretación del rol del HRV en modelos multivariados:} Sugieren una revisión del concepto de que biomarcadores cardiovasculares deben demostrar asociación univariada significativa para ser incluidos en modelos predictivos. La paradoja HRV identificada aporta evidencia sobre la existencia de modificación de efecto y supresión estadística en fisiología del ejercicio, lo cual no había sido claramente demostrado en estudios de vida libre con wearables de consumo.
    
    \item \textbf{Fundamentación del clustering como verdad operativa:} Proporciona evidencia empírica de que clasificaciones derivadas de la estructura inherente de los datos (clustering no supervisado) pueden servir como referencia válida para entrenamiento de modelos supervisados en ausencia de estándares de oro objetivos, particularmente en investigación de vida libre longitudinal.
\end{enumerate}

\subsection{Implicaciones Prácticas}
\label{subsec:implicaciones_practicas}

Los hallazgos tienen aplicaciones prácticas inmediatas en tres ámbitos:

\begin{enumerate}
    \item \textbf{Para profesionales de la salud:} Los resultados sugieren que los datos de wearables de consumo pueden ser incorporados en evaluaciones clínicas de rutina para estratificación de riesgo cardiometabólico. El modelo difuso propuesto puede implementarse en dashboards clínicos que traduzcan datos brutos de actividad física a recomendaciones interpretables (ej. ``Su patrón semanal indica sedentarismo moderado con baja variabilidad cardíaca; considere interrumpir periodos sedentarios cada 30 minutos'').
    
    \item \textbf{Para tomadores de decisiones en salud p├║blica:} Se recomienda la incorporaci├│n de modelos explicables basados en datos de wearables en programas nacionales de vigilancia epidemiol├│gica del sedentarismo, lo cual permite monitoreo poblacional escalable sin requerir equipamiento de investigaci├│n costoso. El enfoque BYOD favorece la equidad al aprovechar infraestructura tecnol├│gica personal existente.
    
    \item \textbf{Para dise├▒adores de intervenciones digitales:} La arquitectura h├¡brida clustering-difuso puede ser adaptada para generar retroalimentaci├│n personalizada en tiempo real, donde las reglas difusas se ajusten din├ímicamente seg├║n el perfil longitudinal del usuario. Esto permitir├¡a que aplicaciones de salud m├│vil evolucionen de recomendaciones gen├®ricas (``camine 10,000 pasos diarios'') hacia intervenciones contextuales (``su patr├│n actual sugiere riesgo elevado; considere actividad ligera durante los pr├│ximos 30 minutos'').
\end{enumerate}

\subsection{Implicaciones Metodol├│gicas}
\label{subsec:implicaciones_metodologicas}

Metodol├│gicamente, este estudio aporta:

\begin{enumerate}
    \item \textbf{Protocolo reproducible para investigación BYOD:} La secuencia de preprocesamiento, derivación de variables, clustering y entrenamiento de sistema difuso está documentada y puede ser replicada en cohortes independientes con dispositivos portátiles alternativos (Garmin, Fitbit), ajustando únicamente las funciones de pertenencia a los centroides observados en la nueva cohorte.
    
    \item \textbf{Estrategia de validaci├│n para cohortes peque├▒as:} Demuestra que LOUO es viable y riguroso en estudios con $N < 15$ participantes pero con alta densidad temporal (m├║ltiples observaciones por individuo), proporcionando m├®tricas realistas de generalizaci├│n que evitan el optimismo de validaciones basadas en partici├│n temporal simple.
    
    \item \textbf{Framework para identificación de modificadores de efecto:} El análisis de ablación combinado con pruebas univariadas ofrece una metodología sistemática para identificar variables que actúan como modificadores de efecto, informando así el diseño de futuros estudios mecanísticos que exploren las interacciones fisiológicas subyacentes.
\end{enumerate}


% ============================================================================
% SECCIÓN 7.6: CUMPLIMIENTO DE OBJETIVOS E HIPÓTESIS
% ============================================================================
\section{Cumplimiento de Objetivos e Hip├│tesis}
\label{sec:cumplimiento}

\subsection{Objetivo General}
\label{subsec:cumplimiento_og}

El objetivo general de \textit{construir y validar un modelo interpretable, basado en l├│gica difusa, para la clasificaci├│n del comportamiento sedentario a partir de datos biom├®tricos de wearables recopilados en condiciones de vida libre} fue alcanzado satisfactoriamente. El sistema de inferencia difusa tipo Mamdani desarrollado demostr├│ capacidad clasificatoria robusta (F1-Score LOOU = 0.780) manteniendo interpretabilidad completa mediante 5 reglas ling├╝├¡sticas auditables.

\subsection{Objetivos Específicos}
\label{subsec:cumplimiento_oe}

\begin{itemize}
    \item \textbf{OE1 - Derivaci├│n de variables semanales:} \textbf{Cumplido.} Se derivaron exitosamente cuatro variables compuestas (Actividad relativa, Super├ívit cal├│rico basal, Delta card├¡aco, HRV-SDNN) mediante normalizaci├│n antropom├®trica y agregaci├│n semanal sobre 9,185 d├¡as de registro, resultando en 1,337 semanas v├ílidas con $<$ 20\% de datos faltantes.
    
    \item \textbf{OE2 - Identificación de perfiles mediante clustering:} \textbf{Cumplido.} El algoritmo K-Means identificó dos perfiles de comportamiento con separación moderada (Silhouette = 0.232) y diferencias estadísticamente significativas en las cuatro variables ($p < 0.001$, tamaños de efecto grandes). Esta clasificación sirvió como verdad operativa para el entrenamiento del sistema difuso.
    
    \item \textbf{OE3 - Diseño del sistema de inferencia difusa:} \textbf{Cumplido.} Se diseñó un sistema Mamdani con 5 reglas lingüísticas fundamentadas fisiológicamente, empleando funciones de pertenencia triangulares calibradas a partir de los percentiles 25, 50 y 75 de cada variable observada en los datos. La estructura \textit{if-then} garantiza interpretabilidad clínica completa.
    
    \item \textbf{OE4 - Evaluaci├│n del desempe├▒o mediante validaci├│n cruzada:} \textbf{Cumplido.} La validaci├│n Leave-One-User-Out demostr├│ concordancia robusta entre el sistema difuso y el clustering (F1 = 0.780, Kappa = 0.56), confirmando que el modelo captur├│ efectivamente los patrones identificados por el m├®todo de agrupamiento.
    
    \item \textbf{OE5 - An├ílisis de sensibilidad mediante ablaci├│n:} \textbf{Cumplido.} El an├ílisis de ablaci├│n sistem├ítico revel├│ que HRV-SDNN, a pesar de su d├®bil asociaci├│n univariada, es cr├¡tico para el desempe├▒o multivariado (ablaci├│n: -51\% F1-Score), caracterizando as├¡ la contribuci├│n relativa de cada variable al modelo predictivo.
\end{itemize}

\subsection{Hip├│tesis}
\label{subsec:cumplimiento_hipotesis}

\textbf{Hipótesis Conceptual (HC):} \textit{La clasificación del comportamiento sedentario generada por el sistema de inferencia difusa, basado en reglas que integran biomarcadores de actividad y función cardiovascular, exhibe una alta concordancia con la clasificación objetiva obtenida mediante un análisis de conglomerados no supervisado sobre los mismos datos.}

\textbf{Resultado:} La hipótesis conceptual fue \textbf{confirmada} con evidencia estadística sólida. El coeficiente Kappa de Cohen entre las clasificaciones del clustering y del sistema difuso alcanzó 0.56 (concordancia moderada-sustancial), significativamente superior al acuerdo esperado por azar ($p < 0.001$). El F1-Score de 0.780 en validación LOUO indica que el modelo difuso logró replicar la estructura clasificatoria identificada por el clustering, generalizando efectivamente a individuos no vistos durante el entrenamiento.

Este resultado confirma la viabilidad de la arquitectura h├¡brida propuesta, donde el conocimiento emp├¡rico descubierto por m├®todos no supervisados puede ser codificado en reglas ling├╝├¡sticas interpretables que preservan la capacidad predictiva. La concordancia observada valida el enfoque de utilizar clustering como verdad operativa en contextos donde est├índares de oro cl├ísicos no est├ín disponibles.


% ============================================================================
% SECCIÓN 7.7: REFLEXIÓN FINAL Y PROYECCIÓN FUTURA
% ============================================================================
\section{Reflexión Final y Líneas Futuras de Investigación}
\label{sec:reflexion_final}

Este estudio demuestra que la convergencia entre inteligencia artificial explicable, dispositivos port├ítiles de consumo y dise├▒os longitudinales de vida libre constituye un eje transformador en la vigilancia de los comportamientos sedentarios. La combinaci├│n de metodolog├¡as BYOD con modelos difusos no solo ampl├¡a la comprensi├│n del fen├│meno desde la perspectiva fisiol├│gica y conductual, sino que tambi├®n abre el camino hacia ecosistemas de salud digital basados en datos de vida real, donde la prevenci├│n y la promoci├│n de la salud se integren de forma continua y personalizada.

La aplicabilidad práctica de estos resultados radica en su potencial para orientar políticas públicas centradas en el uso responsable de la tecnología digital como herramienta de salud. Los modelos explicables, como el propuesto, permiten diseñar dashboards y alertas personalizadas que incentiven la movilidad cotidiana, identifiquen patrones de riesgo y faciliten la toma de decisiones en entornos laborales, educativos y comunitarios. Su naturaleza adaptativa posibilita, además, la personalización de intervenciones conductuales, incorporando recomendaciones dinámicas según la respuesta fisiológica de cada individuo.

En conjunto, los hallazgos aqu├¡ presentados fortalecen la premisa de que la medici├│n objetiva del sedentarismo mediante tecnolog├¡as port├ítiles, interpretada a trav├®s de sistemas inteligentes transparentes, puede convertirse en una herramienta estrat├®gica para la toma de decisiones en salud p├║blica, contribuyendo a reducir la carga de enfermedad asociada a la inactividad y a promover estilos de vida m├ís activos y sostenibles.

\textbf{Líneas futuras de investigación incluyen:}

\begin{enumerate}
    \item \textbf{Validaci├│n prospectiva con est├índares de oro:} Replicar el estudio con mediciones concurrentes de actigraf├¡a de investigaci├│n y calorimetr├¡a indirecta para cuantificar la validez de criterio del modelo difuso versus t├®cnicas de referencia.
    
    \item \textbf{Escalamiento a cohortes heterog├®neas:} Aplicar el protocolo a cohortes con mayor diversidad etaria, socioecon├│mica y cl├¡nica (incluyendo poblaciones con enfermedades cr├│nicas) para evaluar la robustez del modelo en condiciones de alta heterogeneidad.
    
    \item \textbf{Implementación en intervenciones conductuales randomizadas:} Desarrollar una aplicación móvil que utilice el modelo difuso para generar retroalimentación en tiempo real, y evaluar su efectividad en reducir el sedentarismo mediante un ensayo clínico controlado aleatorizado.
    
    \item \textbf{Exploraci├│n de variabilidad temporal:} Analizar patrones circadianos y semanales del comportamiento sedentario mediante descomposici├│n de series temporales, para identificar momentos de mayor susceptibilidad a intervenciones conductuales.
    
    \item \textbf{Integración con biomarcadores inflamatorios y metabólicos:} Vincular los perfiles de sedentarismo identificados con mediciones sanguíneas de proteína C reactiva, glucosa, lípidos y citoquinas inflamatorias, para caracterizar las vías biológicas mediando la asociación sedentarismo-enfermedad.
\end{enumerate}

En s├¡ntesis, los hallazgos aqu├¡ presentados confirman que el monitoreo continuo del sedentarismo mediante dispositivos port├ítiles, interpretado a trav├®s de sistemas de inferencia difusa, representa una v├¡a factible y ├®tica para transformar los datos personales en conocimiento accionable para la salud p├║blica. Este modelo no solo contribuye a la comprensi├│n cient├¡fica del sedentarismo como fen├│meno biopsicosocial, sino que tambi├®n proporciona una base emp├¡rica para dise├▒ar pol├¡ticas p├║blicas orientadas a reducir la inactividad y fomentar entornos digitales saludables. As├¡, la investigaci├│n trasciende el ├ímbito acad├®mico para posicionarse como una herramienta estrat├®gica en la promoci├│n de estilos de vida activos y sostenibles en la era de la salud digital.

% ============================================================================
% FIN DEL CAPÍTULO
% ============================================================================

